% Explicit full-multiplier upgrade of RH.3; genuine tangential comparison.
\subsection{An explicit homology modulus for all multipliers}

\paragraph{LH.1 (Proved): a full-multiplier finite homology detector.}
Let $N=\ker q_{12}$, let $M=N^{\rm ab}\simeq
\widehat{\mathbb Z}^{97}$, and use the genuine normalized
homology action $h_{12}$ of RH.1. Put
\[
                  \boxed{\ n_*=2^{51942}3^{14136}.\ }
\]
The full character $\Delta:\widehat{\mathrm{GT}}_0\to\mathbb F_2$
factors through the finite representation
\[
 h_{12,n_*}:\widehat{\mathrm{GT}}_0\longrightarrow
             \operatorname{GL}_{97}(\mathbb Z/n_*).
\]
Thus the ordered-homology-basis permutation module of RH.3
can be taken at this explicit modulus for all multipliers.
The underlying gentle-stable modular-source quotient
\[
 \Gamma_{n_*}=\widehat B/(N'N^{n_*})
\]
has order
\[
       576n_*^{97}=2^{5038380}3^{1371194}.
\]
These are sufficient bounds. No optimality of this modulus
or quotient order is asserted.

\emph{Proof.}
We first prove the required genuine arithmetic ramification
statement. The reduction $Q/q_{12}(F)\simeq S_3$ implies
$N\subseteq F$. Consequently
\[
                 F/(N'N^{576})
\]
is an actual finite quotient of $F$, of order dividing
$|\Gamma_{576}|=576^{98}$; in particular its only prime
divisors are two and three. The standard tangential Galois
action on $F$ is the literal restriction of the genuine
action used here. Indeed R4, E.65.8 gives
$u\mapsto u^\lambda$ and
$v\mapsto f(u^2,v^2)^{-1}v^\lambda f(u^2,v^2)$; its
restriction to $x=u^2,y=v^2$ is exactly
$x\mapsto x^\lambda$, $y\mapsto f^{-1}y^\lambda f$.
There is no change of outer representative in this comparison.

For every rational prime $p\geq5$, use the smooth pair
$(\mathbb P^1_{\mathbb Z_p},\{0,1,\infty\})$ and the
unit tangential base point $\overrightarrow{01}$.
Petrov's genuine tangential specialization isomorphism
identifies its geometric prime-to-$p$ fundamental group with
that of the special fibre, compatibly with Galois actions,
the latter factoring through the residue Galois group
\cite[Lemma~8.7]{LiftPetrov2023}.
Thus inertia at $p$ acts identically on every finite
$\{2,3\}$ quotient of this tangential $F$. Apply this to
the displayed quotient of $F$. Its subgroup
$N/(N'N^{576})=M/576M$ is therefore fixed pointwise.
It follows that the genuine finite Galois representation
$h_{12,576}|_{G_{\mathbb Q}}$ is unramified outside two
and three. This argument requires neither a multiplier-one
assumption nor any comparison of outer actions on the
modular orbifold.

Let
\[
 K=\ker(h_{12,576}|_{G_{\mathbb Q}}),\qquad
 L=\overline{\mathbb Q}^{\,K},\qquad
 \Lambda=\lambda(K)\subseteq\widehat{\mathbb Z}^{\times},
 \qquad D=|\operatorname{GL}_{97}(\mathbb Z/576)|.
\]
The field $L/\mathbb Q$ is finite Galois, unramified outside
two and three, and $[L:\mathbb Q]$ divides $D$.
Moreover
\[
 [\widehat{\mathbb Z}^{\times}:\Lambda]
       =[L\cap\mathbb Q^{\rm cyc}:\mathbb Q]
       \quad\hbox{divides }D.
\]
The equality is the restriction Galois correspondence; the
divisibility also follows directly because this finite
quotient of the multiplier group is a quotient of
$G_{\mathbb Q}/K$.

For clarity we give the conductor bound in group form.
In the cyclotomic unit product
$\widehat{\mathbb Z}^{\times}=\prod_p\mathbb Z_p^{\times}$,
inertia at a prime $p$ has image the full $p$-factor and
is trivial on the other factors. Since $L$ is unramified
at every $p\geq5$, the closed subgroup $\Lambda$ contains
all these factors, and therefore their product by closedness.
This is also the usual cyclotomic
ramification description of its associated abelian field
\cite[Chapter~I, Theorem~4.16 and its proof]{LiftMilneCFT2020}.
The images of $1+4\mathbb Z_2$ and $1+3\mathbb Z_3$ in
the remaining finite quotient have orders dividing
$2^{v_2(D)}$ and $3^{v_3(D)}$, respectively. The elementary
principal-unit power identities give
\[
 \begin{split}
 (1+4\mathbb Z_2)^{2^{v_2(D)}}
       &=1+2^{v_2(D)+2}\mathbb Z_2\subseteq\Lambda,\\
 (1+3\mathbb Z_3)^{3^{v_3(D)}}
       &=1+3^{v_3(D)+1}\mathbb Z_3\subseteq\Lambda.
 \end{split}
\]
Consequently the principal congruence kernel of modulus
$2^{v_2(D)+2}3^{v_3(D)+1}$ is contained in $\Lambda$.

The standard exact order formula
\[
 |\operatorname{GL}_d(\mathbb Z/p^e)|
   =p^{(e-1)d^2+d(d-1)/2}\prod_{i=1}^d(p^i-1)
\]
and $576=2^6 3^2$ yield
\[
 \begin{split}
 v_2(D)&=5\cdot97^2+\frac{97\cdot96}{2}
               +\sum_{i=1}^{97}v_2(3^i-1)=51940,\\
 v_3(D)&=97^2+\frac{97\cdot96}{2}
               +\sum_{i=1}^{97}v_3(2^i-1)=14135.
 \end{split}
\]
The two cross sums are $239$ and $70$. Exact direct order
valuations and independent principal-power valuations are
checked by \texttt{check\_homology\_bound.py}.
Thus the displayed principal congruence kernel is that of
$n_*$, and $576$ divides $n_*$.

Now suppose $g\in\widehat{\mathrm{GT}}_0$ satisfies
$h_{12,n_*}(g)=1$. The primitive cusp class of $\tau^{12}$
in RH.3 is multiplied by $\lambda(g)$. Its primitivity
therefore gives $\lambda(g)\equiv1\pmod{n_*}$.
By the preceding inclusion $\lambda(g)\in\Lambda$.
Choose an actual $\sigma\in K$ with
$\lambda(\sigma)=\lambda(g)$, as an equality in the full
group $\widehat{\mathbb Z}^{\times}$.
Then $g\sigma^{-1}\in\mathcal H$ and
$h_{12,576}(g\sigma^{-1})=1$.
RH.2 gives $\Delta(g\sigma^{-1})=0$, and R4 gives
$\Delta(\sigma)=0$. Hence $\Delta(g)=0$.
This proves the required finite kernel containment for
every multiplier. The order of $\Gamma_{n_*}$ follows
from its kernel $(\mathbb Z/n_*)^{97}$ and quotient $Q$ of
order $576$.

\paragraph{LH.2 (Proved): the raw homology has no characteristic-zero $\Delta$ constituent.}
For a prime $\ell$, put
$M_\ell=M\otimes_{\widehat{\mathbb Z}}\mathbb Z_\ell$
and $M_{\mathbb Q_\ell}=M_\ell\otimes_{\mathbb Z_\ell}\mathbb Q_\ell$.
Then
\[
                 M_{\mathbb Q_\ell}^{G_{\mathbb Q}}=0.
\]
More strongly, no subquotient of this finite-dimensional
$G_{\mathbb Q}$ representation is the trivial representation.
Consequently the raw homology $M_{\mathbb Q_\ell}$ has no
$\widehat{\mathrm{GT}}_0$ subquotient affording the character
$(-1)^\Delta$, and in particular no such line. There is also
no free rank-one $\mathbb Z_\ell$ submodule of $M_\ell$
with this action. These statements concern the full gentle
group and its full Galois subgroup; they do not assert the
corresponding statement after restriction to $\mathcal H$.

\emph{Proof.}
Let $X=\mathbb P^1_{\mathbb Q}\setminus\{0,1,\infty\}$,
with its standard rational tangential base point. The genuine
action comparison in LH.1 identifies its arithmetic
fundamental group with $F\rtimes G_{\mathbb Q}$, using
exactly the action under consideration. Since $N\subseteq F$
is stable under this action, the open subgroup
\[
                  N\rtimes G_{\mathbb Q}
                  \ \subseteq\ F\rtimes G_{\mathbb Q}
\]
corresponds to a finite \emph{geometrically connected} etale
cover $Y\to X$ over $\mathbb Q$, of degree $[F:N]=96$.
Indeed it projects onto $G_{\mathbb Q}$, and its geometric
intersection is exactly $N$. This is the usual finite-cover
correspondence, applied to the specified genuine section
\cite[\S3]{LiftMilneLEC2013}. Its geometric first homology
is $M_{\mathbb Q_\ell}$ with precisely the given Galois
action. Changing a geometric base point only conjugates
by $N$, which has no effect on its abelianization.

Let $\overline Y$ be the smooth projective compactification,
let $D=\overline Y\setminus Y$, and let
$J=\operatorname{Jac}(\overline Y)$. The Gysin sequence
for this smooth curve and its finite boundary, dualized,
gives the natural Galois-equivariant exact sequence
\[
 0\longrightarrow
 \mathbb Q_\ell(1)\otimes
 \left(\mathbb Q_\ell[D(\overline{\mathbb Q})]
       /\mathbb Q_\ell\!\cdot\!\sum_{d\in D(\overline{\mathbb Q})}d\right)
 \longrightarrow M_{\mathbb Q_\ell}
 \longrightarrow V_\ell J\longrightarrow0.
\]
Here $V_\ell J=T_\ell J\otimes\mathbb Q_\ell$; the compact
curve comparison follows from its Jacobian's universal
abelian covering property and the Tate-module description
of first homology
\cite[Corollary~16.2]{LiftMilneLEC2013}
\cite[Proposition~9.1 and Corollary~9.6]{LiftMilneJac2021}
\cite[Theorem~15.1 and Remark~15.5]{LiftMilneAV2022}.

Choose a rational prime $p\ne\ell$ of good reduction for
$\overline Y$ at which its finite boundary is unramified.
Such primes exist because these data spread out over
$\mathbb Z[1/S]$ for a finite set $S$. Let $\mathrm{Fr}_p$
be an arithmetic Frobenius lift. On the boundary term its
eigenvalues are $p$ times roots of unity: the finite point set
is permuted, and the Tate twist is multiplication by $p$.
Thus this term has no eigenvalue one.

For the compact term, good reduction identifies the
$\ell$-power torsion with that of the special-fibre Jacobian
$J_p/\mathbb F_p$, compatibly with arithmetic Frobenius.
One can see this directly from the finite etale group schemes
$J[\ell^m]$ over $\mathbb Z_p$
\cite[Proposition~20.7]{LiftMilneAV2022}.
The endomorphism $\operatorname{Frob}_p-1$ of $J_p$ is an
isogeny: its geometric kernel is the finite set
$J_p(\mathbb F_p)$, and its differential is $-1$, so this
kernel is a finite etale group scheme. An endomorphism of an
abelian variety with finite kernel is surjective and finite.
It therefore becomes invertible on $V_\ell J_p$.
Equivalently, arithmetic Frobenius has no eigenvalue one
on $V_\ell J$.

The characteristic polynomial on the middle term is the
product of those on the two end terms. Hence
$\det(\mathrm{Fr}_p-1\mid M_{\mathbb Q_\ell})\ne0$.
Every invariant subquotient has a characteristic polynomial
dividing this polynomial, so no trivial Galois subquotient
can occur. R4 says $\Delta|_{G_{\mathbb Q}}=0$; a full
gentle-group $\Delta$ constituent would therefore restrict
to a trivial Galois constituent, a contradiction. A free
$\mathbb Z_\ell$ line would yield such a line after tensoring
with $\mathbb Q_\ell$.

This does not exclude a constituent with finite coefficients:
the nonzero determinant above need not be a unit modulo a
chosen integer. It is also consistent with LH.1, whose
sign constituent lies in the permutation module of finite
ordered homology bases, rather than in the raw homology.
